\subsection*{S22. Name mangling and collisions}

Tables are named after the field whose value they hold, so two parents using
the same member name want the same table:

\begin{lstlisting}[style=json]
{ "user":    { "profile": { "name": "ada" } },
  "company": { "profile": { "industry": "rail" } } }
\end{lstlisting}

\noindent
Two tables named \texttt{profile}, with different columns.

\options
\begin{enumerate}[label=\textbf{(\alph*)}]
  \item \textbf{Qualify by path} --- \texttt{user\_profile},
        \texttt{company\_profile}. Always unambiguous; names grow with nesting
        depth. \selected
  \item \textbf{Suffix on collision} --- \texttt{profile},
        \texttt{profile\_2}. Short, but the numbering depends on traversal
        order, so adding an unrelated field can renumber existing modules.
  \item \textbf{Merge structurally identical tables}, qualifying only when
        they differ.
  \item \textbf{Reject on collision}, requiring an annotation to
        disambiguate.
\end{enumerate}

\decision{Qualify by path. A table's name is its path from the document root,
with each segment mangled by S11's scheme and the segments joined by
\texttt{\_}. The module name is that name with its first character
capitalized.}

\begin{lstlisting}[style=json]
.                        -> table root            module Root
.user.profile            -> table user_profile    module User_profile
.company.profile         -> table company_profile module Company_profile
.order.items[]           -> table order_items     module Order_items
\end{lstlisting}

Array brackets contribute no segment: the element table is named for the field
holding the array. No ambiguity arises from this, because a path cannot be
both an array and an object (S18 would reject it), and arrays of arrays are
rejected outright (S17). Where S19 folds two paths together, the table takes
the name of the ancestor path, which is the one that survives the
unification.

\paragraph{Collisions cannot occur.}
Under (a) the scheme is injective by construction, so option (d) never fires.
Three earlier decisions compose to give this:

\begin{itemize}
  \item Paths are distinct by definition, so distinct tables have distinct
        paths.
  \item S11's mangling is injective on each segment.
  \item S11 escapes \texttt{\_} itself, as \texttt{\_5f}. A literal underscore
        in a member name therefore cannot be confused with the separator
        joining two segments: the single field \texttt{"a\_b"} mangles to
        \texttt{a\_5fb}, while the path \texttt{.a.b} joins to
        \texttt{a\_b}.
\end{itemize}

\paragraph{Capitalization is injective too.}
Module names must begin with an uppercase letter, and naive capitalization is
not injective --- the distinct members \texttt{b} and \texttt{B} would both
give module \texttt{B}. S11 already resolves this before it arises. Its second
rule prefixes \texttt{f\_} to anything not beginning with a lowercase letter,
so \texttt{B} mangles to \texttt{f\_B} while \texttt{b} mangles to \texttt{b}.
Every mangled name therefore begins with a lowercase ASCII letter, and
capitalizing that one character is a bijection on the set of possible names:

\begin{lstlisting}[style=json]
member "b"  ->  b     -> module B
member "B"  ->  f_B   -> module F_B
\end{lstlisting}

\rationale{Path qualification is purely structural. A table's name depends
only on where it sits in the schema, not on what else the corpus happens to
contain, which is the property protected in S9, S19 and S20. Option (b) fails
it --- the numbering depends on traversal order, so adding a field elsewhere
can rename an existing module and break code written against it --- and option
(c) fails it more seriously, since whether two tables merge would depend on how
much of the corpus had been seen.}

\limitation{Names grow with nesting depth, and the Delivery note establishes
that these modules are read by people in an editor rather than consumed as a
build artifact. A field four levels down becomes
\texttt{User\_profile\_address\_geo}, which is accurate and unpleasant.

Option (c) produces markedly nicer schemas and remains available, but adopting
it would mean accepting that a table's name depends on the rest of the corpus.
A middle path, not adopted here, would be to allow the configuration file to
supply an explicit name for a path, in the same manner as the map and
recursion annotations of S9 and S19 --- keeping generation structural while
letting a person shorten the names that matter.}
